%% toc: P15 166 166.3
% Proofs for continuous entropy relative to a reference measure, Shannon's Section 20 properties and coordinate changes.
\section{Continuous entropy and coordinates: proofs}
\label{sec:annex-p15-w02b-p15-042-continuous-entropy}

Use the definitions of
Chapter~\ref{sec:p15-w02b-p15-042-continuous-entropy}; all logarithms are
natural and \(\eta(t)=-t\log t\), \(\eta(0)=0\). Two facts about extended
integrals are used repeatedly. Call a measurable \(G\) \emph{quasi-integrable}
when at least one of \(\int G^+\) and \(\int G^-\) is finite.

\emph{(E1) Subtracting a nonnegative function.} If \(F\) is integrable and
\(\psi\geq0\) is measurable, then \(G=F-\psi\) is quasi-integrable and
\(\int G=\int F-\int\psi\in[-\infty,\infty)\). Indeed \(G^+\leq F^+\). If
\(\int\psi<\infty\), linearity of the integral for integrable functions
applies. If \(\int\psi=\infty\), then \(\psi\leq G^-+F^+\) pointwise, so
\(\int G^-=\infty\) and both sides equal \(-\infty\).

\emph{(E2) Adding an integrable function.} If \(G\) is quasi-integrable and
\(K\) is integrable, then \(G+K\) is quasi-integrable and
\(\int(G+K)=\int G+\int K\). If \(G\) is integrable this is linearity.
If \(\int G^-<\infty=\int G^+\), then \((G+K)^-\leq G^-+K^-\) is integrable,
and \(G^+\leq(G+K)^++K^-\) forces \(\int(G+K)^+=\infty\); both sides are
\(+\infty\). The case \(\int G^+<\infty=\int G^-\) is the same argument
applied to \(-G\) and \(-K\). By (E2) with the roles exchanged, the sum of an
integrable function and a quasi-integrable one is quasi-integrable with the
additive integral.

\emph{(T) The tangent inequality.} For \(s>0\) and \(t\geq0\) put
\(\varphi(t,s)=t\log(t/s)-t+s\), with \(0\log0=0\). For fixed \(s\), the
function \(t\mapsto\varphi(t,s)\) is continuous on \([0,\infty)\) and has
derivative \(\log(t/s)\) on \((0,\infty)\), negative on \((0,s)\) and
positive on \((s,\infty)\). By the mean value theorem it decreases on
\([0,s]\) and increases on \([s,\infty)\), so its minimum is
\(\varphi(s,s)=0\), attained only at \(t=s\). Thus \(\varphi\geq0\), with
equality exactly when \(t=s\). Equivalently,
\(\eta(t)\leq-t\log s-t+s\), with equality exactly when \(t=s\).

\begin{annexitem}{annex:p15-w02b-p15-042-gibbs}{thm:p15-w02b-p15-042-gibbs}{Comparison with a subprobability density}
\begin{proof}
Let \(N=\{p>0,\ q=0\}\). On \(N\) the product \(p\log q\) equals \(-\infty\);
since \(p\log q\) is integrable, \(\mu(N)=0\). Define \(\psi=\varphi(p,q)\)
on \(\{q>0\}\) and \(\psi=0\) on \(\{q=0\}\); \(\psi\geq0\) by (T). On
\(\{q>0\}\),
\[
 \eta(p)=-p\log q-p\log\frac pq=-p\log q-\psi+q-p ,
\]
both when \(p>0\) and, with \(0\log0=0\), when \(p=0\). On
\(\{q=0\}\setminus N\) we have \(p=0\), and both sides vanish. Hence
\(\eta(p)=F-\psi\) \(\mu\)-a.e., where \(F=-p\log q+q-p\) satisfies
\(|F|\leq|p\log q|+q+p\) and is integrable. By (E1), \(H_\mu(p)\) is defined
and
\[
 H_\mu(p)=-\int p\log q\,d\mu+\Bigl(\int q\,d\mu-1\Bigr)-\int\psi\,d\mu
         \leq-\int p\log q\,d\mu<\infty ,
\]
because \(\int q\leq1\) and \(\int\psi\geq0\).

Equality holds exactly when \(\int\psi=0\) and \(\int q=1\). The first
condition means \(\psi=0\) a.e., that is \(p=q\) a.e.\ on \(\{q>0\}\) by (T).
Then \(\int_{\{q>0\}}q=\int_{\{q>0\}}p=1\), because \(p=0\) a.e.\ on
\(\{q=0\}\); so \(\int q=1\) holds automatically and \(p=q\) a.e.\ on all of
\(S\). Conversely, \(p=q\) a.e.\ gives \(\psi=0\) a.e.\ and \(\int q=1\).
\end{proof}
\end{annexitem}

\begin{annexitem}{annex:p15-w02b-p15-042-reference-change}{thm:p15-w02b-p15-042-reference-change}{Changing the reference measure}
\begin{proof}
The measure \(\nu\) is sigma-finite: if \(S=\bigcup_mS_m\) with
\(\mu(S_m)<\infty\), the sets \(S_m\cap\{k-1<f\leq k\}\) cover \(S\) and have
\(\nu\)-measure at most \(k\mu(S_m)\). By definition \(f\) is a derivative of
\(\nu\) with respect to \(\mu\), so item~1 of \cref{thm:p15-density-calculus}
gives \(\int h\,d\nu=\int hf\,d\mu\) for measurable \(h\geq0\). With
\(h=\mathbf 1_A\,p/f\), this yields \(\int_A(p/f)\,d\nu=\int_Ap\,d\mu=P(A)\),
so \(p/f\) is a \(\nu\)-density of \(P\).

With \(h=\eta(p/f)^\pm\) and \(f>0\),
\[
 \int\eta(p/f)^\pm\,d\nu=\int\bigl(f\eta(p/f)\bigr)^\pm\,d\mu .
\]
Where \(p>0\), \(f\eta(p/f)=-p\log p+p\log f=\eta(p)+p\log f\); where
\(p=0\) both sides are zero. So \(H_\nu(P)\) is defined exactly when
\(\eta(p)+p\log f\) is quasi-integrable relative to \(\mu\), and it then
equals \(\int(\eta(p)+p\log f)\,d\mu\). Under the first hypothesis
\(\eta(p)\) is quasi-integrable and \(p\log f\) is integrable; under the
second, \(\eta(p)\) is integrable and \(p\log f\) is quasi-integrable. In
both cases (E2) gives quasi-integrability and
\(H_\nu(P)=H_\mu(P)+\int p\log f\,d\mu\).
\end{proof}
\end{annexitem}

\begin{annexitem}{annex:p15-w02b-p15-042-defined-cases}{thm:p15-w02b-p15-042-defined-cases}{Finite, infinite and undefined entropies}
\begin{proof}
\emph{Item 1.} The function \(\eta\) is continuous on \([0,\infty)\), with
derivative \(-\log t-1\) on \((0,\infty)\), so its maximum is
\(\eta(1/e)=1/e\). Hence \(\eta(p)^+\leq1/e\) and
\(\int\eta(p)^+\,d\mu\leq\mu(S)/e<\infty\).

\emph{An integral used below.} For \(b>0\), \(\int_{(b,\infty)}u^{-2}\,du=1/b\)
and \(\int_{(b,\infty)}u^{-1}\,du=\infty\): on \([b,N]\) the integrands are
continuous, the integrals equal \(1/b-1/N\) and \(\log(N/b)\) by the
fundamental theorem of calculus, and monotone convergence
(\cref{thm:p15-monotone-integral}) lets \(N\to\infty\).

\emph{Item 2.} The map \(T(x)=-\log x\) is a decreasing \(C^1\) bijection
from \((0,\tfrac12)\) onto \((\log2,\infty)\) with \(|T'(x)|=1/x\). By the
substitution formula of \cref{thm:p15-criteria-real-integration}, for Borel
\(h\geq0\),
\[
 \int_{(0,1/2)}h(-\log x)\,\frac{dx}{x}=\int_{(\log2,\infty)}h(u)\,du .
\]
Since \((\log x)^2=u^2\) for \(u=-\log x\), \(p_1(x)=h(-\log x)/x\) with
\(h(u)=\log2/u^2\); hence \(\int p_1=\log2\cdot(1/\log2)=1\). For
\(0<x<\tfrac12\) and \(u=-\log x\),
\(\log p_1(x)=\log\log2+u-2\log u\), so \(\eta(p_1(x))=p_1(x)k(u)\) with
\(k(u)=-u+2\log u-\log\log2\). Since \(k'(u)=-1+2/u\), the maximum of
\(k\) is \(k(2)=-2+2\log2-\log\log2\approx-0.247<0\). So
\(\eta(p_1)\leq0\), the positive part vanishes, and \(H_\lambda(p_1)\) is
defined. For \(u\geq16\),
\(-k(u)-u/2=u/2-2\log u+\log\log2\) is increasing (its derivative is
\(\tfrac12-2/u>0\)) and is positive at \(u=16\)
(\(8-2\log16+\log\log2\approx8-5.545-0.367\)). Applying the substitution
to \(h=-\log2\cdot k(u)/u^2\geq0\),
\[
 \int\eta(p_1)^-\,dx=\log2\int_{(\log2,\infty)}\frac{-k(u)}{u^2}\,du
   \geq\frac{\log2}2\int_{(16,\infty)}\frac{du}u=\infty .
\]
Thus \(H_\lambda(p_1)=-\infty\).

\emph{Item 3.} The map \(T(x)=\log x\) is an increasing \(C^1\) bijection
from \((e,\infty)\) onto \((1,\infty)\) with \(T'(x)=1/x\), and
\(p_2(x)=h(\log x)/x\) for \(h(u)=u^{-2}\). The substitution gives
\(\int p_2=1\). For \(u=\log x>1\), \(\log p_2(x)=-u-2\log u\), so
\(\eta(p_2(x))=p_2(x)(u+2\log u)\geq0\), and
\[
 \int\eta(p_2)\,dx=\int_{(1,\infty)}\frac{u+2\log u}{u^2}\,du
   \geq\int_{(1,\infty)}\frac{du}u=\infty .
\]
Thus \(H_\lambda(p_2)=+\infty\).

\emph{Item 4.} The supports are disjoint, so \(\int p_3=1\). On
\((0,\tfrac12)\), \(p_3=p_1/2\) and
\(\eta(p_3)=\tfrac12\eta(p_1)+\tfrac{\log2}2p_1\); hence
\(\eta(p_3)^-+\tfrac{\log2}2p_1\geq\tfrac12\eta(p_1)^-\), and integration
gives \(\int\eta(p_3)^-=\infty\) by item~2. On \((e,\infty)\),
\(\eta(p_3)=\tfrac12\eta(p_2)+\tfrac{\log2}2p_2\geq\tfrac12\eta(p_2)\geq0\),
so \(\int\eta(p_3)^+=\infty\) by item~3.

\emph{Item 5.} Under a law with a density relative to \(\lambda\) every
singleton has probability zero, and \(\delta_c\) has no such density
(\cref{thm:p15-density-integrals}); so no law with an atom has one. The Cantor
law has no Lebesgue density by \cref{thm:p15-cantor-singular-law}.
Definition~\ref{def:p15-w02b-p15-042-reference-entropy} leaves the entropy
undefined in these cases.
\end{proof}
\end{annexitem}

\begin{annexitem}{annex:p15-w02b-p15-042-volume}{thm:p15-w02b-p15-042-volume}{Item 1: the bound by the logarithm of the volume}
\begin{proof}
Since \(\int_{S\setminus A}p\,d\mu=P(S\setminus A)=0\), \(p=0\) a.e.\ off
\(A\). Let \(q=\mathbf 1_A/v\), so \(\int q=1\). On \(A\),
\(p\log q=-p\log v\); off \(A\), \(p\log q=0\) a.e.\ by the convention
\(0\log0=0\). Thus \(p\log q\) is integrable with
\(-\int p\log q\,d\mu=\log v\). \Cref{thm:p15-w02b-p15-042-gibbs} gives that
\(H_\mu(P)\) is defined, \(H_\mu(P)\leq\log v\), and equality exactly when
\(p=q\) a.e. The value \(-\infty\) occurs for
\(\mu=\lambda\), \(A=(0,\tfrac12)\) and the density \(p_1\) of
\cref{thm:p15-w02b-p15-042-defined-cases}.
\end{proof}
\end{annexitem}

\begin{annexitem}{annex:p15-w02b-p15-042-joint}{thm:p15-w02b-p15-042-joint}{Item 2: subadditivity and its equality case}
\begin{proof}
Write \(\rho=\mu_1\times\mu_2\). By Tonelli (\cref{thm:p15-tonelli}) the
marginal functions are measurable, and
\(P(x\in A)=\int_A\int p(x,y)\,d\mu_2(y)\,d\mu_1(x)\), so \(p(x)\) is a
density of the law of \(x\); it is finite a.e.\ because its integral is
\(1\), and its values on the null set where it is infinite may be set to
zero. The same holds for \(p(y)\).

Let \(q(x,y)=p(x)p(y)\); by Tonelli \(\int q\,d\rho=1\). The set
\(\{p(x)=0\}\times S_2\) has \(P\)-probability
\(\int_{\{p(x)=0\}}p(x)\,d\mu_1=0\), so \(p(x,y)=0\) for \(\rho\)-almost every
point of it; similarly for \(S_1\times\{p(y)=0\}\). Hence, \(\rho\)-a.e.,
either \(p(x,y)=0\) or \(p(x),p(y)>0\), and
\[
 p(x,y)\log q(x,y)=p(x,y)\log p(x)+p(x,y)\log p(y)\quad\rho\text{-a.e.},
\]
with \(0\log0=0\). By Tonelli,
\(\int p(x,y)|\log p(x)|\,d\rho=\int p(x)|\log p(x)|\,d\mu_1=\int|\eta(p(x))|\,d\mu_1\),
which is finite because \(H(x)\) is finite; similarly for \(y\). So
\(p\log q\) is integrable and, by Fubini (\cref{thm:p15-fubini}),
\(-\int p\log q\,d\rho=H(x)+H(y)\).
\Cref{thm:p15-w02b-p15-042-gibbs} now gives that \(H(x,y)\) is defined, lies in
\([-\infty,\infty)\), and is at most \(H(x)+H(y)\), with equality exactly when
\(p=q\) \(\rho\)-a.e.

If \(p=q\) \(\rho\)-a.e., Tonelli gives
\(P(x\in A,\ y\in B)=\int_Ap(x)\,d\mu_1\int_Bp(y)\,d\mu_2
 =P(x\in A)P(y\in B)\), which is independence. Conversely, under independence
\(P\) and the measure \(q\rho\) are probability measures agreeing on every
rectangle \(A\times B\); by the uniqueness in
\cref{thm:p15-product-existence} both equal the product of the marginal
laws. Then \(p\) and \(q\) are densities of the same measure relative to the
sigma-finite \(\rho\), and they agree \(\rho\)-a.e.\ by the uniqueness part of
\cref{thm:p15-radon-nikodym}.
\end{proof}
\end{annexitem}

\begin{annexitem}{annex:p15-w02b-p15-042-averaging}{thm:p15-w02b-p15-042-averaging}{Item 3: generalized averaging does not decrease entropy}
\begin{proof}
By Tonelli, \(p'\) is measurable and
\[
 \int p'\,d\mu=\int p(x)\int a(x,y)\,d\mu(y)\,d\mu(x)=\int p\,d\mu=1;
\]
\(p'\) is finite a.e., and its values on the infinite-value null set may be
set to zero. So \(p'\) is a \(\mu\)-density.

Again by Tonelli,
\(\iint a(x,y)|\eta(p(x))|\,d\mu(x)\,d\mu(y)=\int|\eta(p)|\,d\mu<\infty\).
Hence for a.e.\ \(y\) the function \(x\mapsto a(x,y)\eta(p(x))\) is
integrable; put \(g(y)=\int a(x,y)\eta(p(x))\,d\mu(x)\) for such \(y\) and
\(g(y)=0\) otherwise. By Fubini (\cref{thm:p15-fubini}), \(g\) is integrable
and \(\int g\,d\mu=\int\eta(p)\,d\mu=H_\mu(p)\).

Fix such a \(y\) with \(p'(y)<\infty\). If \(s=p'(y)>0\), (T) gives
\(\eta(p(x))\leq-p(x)\log s-p(x)+s\) for every \(x\). Multiply by
\(a(x,y)\geq0\) and integrate in \(x\); each term on the right is integrable,
and \(\int a(x,y)p(x)\,d\mu(x)=s\), \(\int a(x,y)\,d\mu(x)=1\). Thus
\[
 g(y)\leq-s\log s-s+s=\eta(p'(y)).
\]
If \(p'(y)=0\), then \(a(x,y)p(x)=0\) for a.e.\ \(x\); at such \(x\),
\(a(x,y)\eta(p(x))=0\) (either \(a(x,y)=0\) or \(p(x)=0\)), so
\(g(y)=0=\eta(p'(y))\). Hence \(\eta(p')\geq g\) a.e. Then
\(\eta(p')^-\leq g^-\) is integrable, so \(H_\mu(p')\) is defined and
greater than \(-\infty\), and by (E2) applied to
\(\eta(p')=g+(\eta(p')-g)\) with \(\eta(p')-g\geq0\),
\(H_\mu(p')=\int g+\int(\eta(p')-g)\geq H_\mu(p)\).
\end{proof}
\end{annexitem}

\begin{annexitem}{annex:p15-w02b-p15-042-conditional}{thm:p15-w02b-p15-042-conditional}{Item 4: joint and conditional entropies}
\begin{proof}
As in the proof of \cref{thm:p15-w02b-p15-042-joint}, \(p(x)>0\) for
\(\rho\)-almost every point with \(p(x,y)>0\), and \(p(x,y)\log p(x)\) is
integrable with integral \(\int p(x)\log p(x)\,d\mu_1=-H(x)\). Where
\(p(x,y)>0\),
\[
 -p(x,y)\log\frac{p(x,y)}{p(x)}=\eta(p(x,y))+p(x,y)\log p(x);
\]
where \(p(x,y)=0\) both sides vanish. Since \(H(x,y)\) is finite, \(\eta(p(x,y))\)
is integrable, so the integrand of \(H_x(y)\) is integrable and
\(H_x(y)=H(x,y)-H(x)\), a finite number. Exchanging the roles of \(x\) and
\(y\) gives \(H_y(x)=H(x,y)-H(y)\). By \cref{thm:p15-w02b-p15-042-joint},
\(H_x(y)=H(x,y)-H(x)\leq H(y)\), with equality exactly under independence.

For the example, \(p(x,y)=2\) on \(R=(0,1)\times(0,\tfrac12)\) and \(0\)
elsewhere; \(\lambda^2(R)=\tfrac12\) by \cref{thm:p15-finite-product-measures}.
The marginal densities are \(p(x)=1\) on \((0,1)\) and \(p(y)=2\) on
\((0,\tfrac12)\). By the equality case of \cref{thm:p15-w02b-p15-042-volume},
\(H(x,y)=\log\tfrac12\), \(H(x)=0\) and \(H(y)=\log\tfrac12\), all finite.
Hence \(H_x(y)=-\log2\).
\end{proof}
\end{annexitem}

\begin{annexitem}{annex:p15-w02b-p15-042-gaussian}{thm:p15-w02b-p15-042-gaussian}{Items 5 and 6: Gaussian maximizers and Gaussian entropies}
\begin{proof}
\emph{Item 1.} Let \(q=\phi_{c,\sigma^2}\), a probability density with mean
\(c\) and variance \(\sigma^2\) by \cref{thm:p15-normal-normalization}. Then
\(\log q(x)=-\log(\sqrt{2\pi}\,\sigma)-(x-c)^2/(2\sigma^2)\), so
\(|p\log q|\leq p\,|\log(\sqrt{2\pi}\sigma)|+p(x-c)^2/(2\sigma^2)\) is
integrable and
\[
 -\int p\log q\,dx=\log(\sqrt{2\pi}\,\sigma)+\tfrac12
                  =\log\sqrt{2\pi e}\,\sigma .
\]
\Cref{thm:p15-w02b-p15-042-gibbs} gives the defined entropy, the bound and the
equality condition; \(q\) itself satisfies the constraint, so the bound is
attained.

\emph{A factorization.} Every symmetric positive definite \(n\times n\)
matrix \(a\) equals \(B^{\mathsf T}B\) for an invertible \(B\). For \(n=1\)
take \(B=(\sqrt{a_{11}})\). For \(n>1\), \(a_{11}=e_1^{\mathsf T}ae_1>0\).
Write \(x=(x_1,x')\) and let \(a'\) be the symmetric
\((n-1)\times(n-1)\) matrix
\(a'_{ij}=a_{ij}-a_{1i}a_{1j}/a_{11}\) (\(i,j\geq2\)). Expanding the square
and using \(a_{i1}=a_{1i}\),
\[
 \sum_{i,j}a_{ij}x_ix_j
   =a_{11}\Bigl(x_1+\sum_{j\geq2}\frac{a_{1j}}{a_{11}}x_j\Bigr)^2
    +\sum_{i,j\geq2}a'_{ij}x_ix_j .
\]
For \(x'\neq0\), choosing \(x_1=-\sum_{j\geq2}a_{1j}x_j/a_{11}\) makes
\(x\neq0\) and shows \(\sum a'_{ij}x_ix_j>0\); so \(a'\) is positive
definite and, by induction, \(a'=B'^{\mathsf T}B'\) with \(B'\) invertible.
Let \(B\) have first row
\((\sqrt{a_{11}},a_{12}/\sqrt{a_{11}},\ldots,a_{1n}/\sqrt{a_{11}})\) and
remaining rows \((0\mid B')\). Then \(|Bx|^2\) equals the displayed right
side, so \(x^{\mathsf T}(B^{\mathsf T}B-a)x=0\) for all \(x\). A symmetric
matrix \(S\) with this property is zero: \(S_{ii}=e_i^{\mathsf T}Se_i=0\) and
\((e_i+e_j)^{\mathsf T}S(e_i+e_j)=2S_{ij}=0\). Finally
\(\det B=\sqrt{a_{11}}\det B'\neq0\).

\emph{Item 2.} Let \(a=B^{\mathsf T}B\) as above, so
\(|a_{ij}|=(\det B)^2\) and \(\sum a_{ij}x_ix_j=|Bx|^2\). Let
\(g(z)=(2\pi)^{-n/2}e^{-|z|^2/2}=\prod_k\phi_{0,1}(z_k)\); by Tonelli and
\cref{thm:p15-normal-normalization}, \(\int g\,dz=1\). Since
\(q_a(x)=|\det B|\,g(Bx)\), \cref{thm:p15-w02b-p15-042-linear-lebesgue} gives
\(\int q_a=|\det B|\cdot|\det B|^{-1}\int g=1\). For Borel \(h\geq0\), the same
lemma applied to \(G(z)=h(B^{-1}z)g(z)\) gives
\(\int h(x)q_a(x)\,dx=\int h(B^{-1}z)g(z)\,dz\); for signed \(h\) apply this to
\(h^\pm\). Write \(C=B^{-1}\). Then
\(x_ix_j\) at \(x=Cz\) is \(\sum_{k,l}C_{ik}C_{jl}z_kz_l\), and
\(|z_kz_l|\leq(z_k^2+z_l^2)/2\) is \(g\)-integrable. By Tonelli, Fubini and
\cref{thm:p15-normal-normalization}, \(\int z_kz_lg\,dz\) is
\(\mathbb EZ^2=1\) for \(k=l\) and \((\mathbb EZ)^2=0\) for \(k\neq l\), where
\(Z\) is standard normal. Hence
\[
 \int x_ix_jq_a(x)\,dx=\sum_kC_{ik}C_{jk}=(CC^{\mathsf T})_{ij}
   =\bigl((B^{\mathsf T}B)^{-1}\bigr)_{ij}=A_{ij}.
\]

\emph{Item 3.} Here
\(\log q_a(x)=-\tfrac n2\log2\pi+\tfrac12\log|a_{ij}|
 -\tfrac12\sum a_{ij}x_ix_j\), and
\(|\sum a_{ij}x_ix_j|\leq\sum_{i,j}|a_{ij}|(x_i^2+x_j^2)/2\) is
\(p\)-integrable because \(\int x_i^2p=A_{ii}<\infty\). So \(p\log q_a\) is
integrable and, using \(\sum_ja_{ij}A_{ji}=1\) for each \(i\) and
\(A_{ij}=A_{ji}\),
\[
 -\int p\log q_a\,dx=\tfrac n2\log2\pi-\tfrac12\log|a_{ij}|
    +\tfrac12\sum_{i,j}a_{ij}A_{ij}
   =\log\,(2\pi e)^{n/2}|a_{ij}|^{-1/2}.
\]
\Cref{thm:p15-w02b-p15-042-gibbs} gives the bound and the equality condition,
and item~2 shows that \(q_a\) satisfies the constraints.

\emph{Item 4.} If \((A_{ij})\) is symmetric and not positive definite, there
is \(v\neq0\) with \(\sum v_iv_jA_{ij}\leq0\). For a density \(p\) with
these moments, \(\int(v^{\mathsf T}x)^2p\,dx=\sum v_iv_jA_{ij}\leq0\), so
\(p=0\) a.e.\ off the hyperplane \(L=\{v^{\mathsf T}x=0\}\). Choose \(k\) with
\(v_k\neq0\). Using \cref{thm:p15-finite-product-measures} to write
\(\lambda^n\) as the product of \(\lambda\) in coordinate \(k\) and
\(\lambda^{n-1}\) in the others, each section of \(L\) in coordinate \(k\) is
one point, so Tonelli gives \(\lambda^n(L)=0\). Then \(\int p=0\), a
contradiction.

\emph{Item 5.} Apply the computations of items 1 and 3 with \(p=q\) and
\(p=q_a\), which satisfy the constraints; the integrability shown there makes
\(H(q)\) and \(H(q_a)\) finite and equal to the displayed values.
\end{proof}
\end{annexitem}

\begin{annexitem}{annex:p15-w02b-p15-042-exponential}{thm:p15-w02b-p15-042-exponential}{Item 7: the exponential maximizer on a half line}
\begin{proof}
Let \(q(x)=a^{-1}e^{-x/a}\) for \(x>0\) and \(q(x)=0\) otherwise. This is the
gamma density of Definition~\ref{def:p15-gamma-density} with shape \(1\) and
rate \(1/a\); by \cref{thm:p15-gamma-normalization} it integrates to \(1\)
and has mean \(a\). For \(x>0\), \(p\log q=p(-\log a-x/a)\), which is
integrable because \(\int_0^\infty xp\,dx=a\); for \(x\leq0\), \(p=0\) a.e.\
and \(p\log q=0\) a.e. Thus \(-\int p\log q\,dx=\log a+1=\log ea\), and
\cref{thm:p15-w02b-p15-042-gibbs} gives the bound and the equality condition;
\(q\) satisfies the constraints.
\end{proof}
\end{annexitem}

\begin{annexitem}{annex:p15-w02b-p15-042-linear-lebesgue}{thm:p15-w02b-p15-042-linear-lebesgue}{Linear and product coordinate changes}
\begin{proof}
Say that an invertible \(M\) is \emph{good} if
\(\int g(Mx)\,dx=|\det M|^{-1}\int g\,dz\) for all Borel \(g\geq0\); every
\(g\circ M\) is Borel because \(M\) is continuous.

\emph{Elementary matrices are good.} By
\cref{thm:p15-finite-product-measures}, permuting coordinates carries
\(\lambda^n\) to itself, so a matrix that exchanges two coordinates
(determinant \(\pm1\)) is good; the same theorem lets us regard \(\lambda^n\)
as the product of \(\lambda\) in a chosen coordinate \(k\) and
\(\lambda^{n-1}\) in the others, so that Tonelli (\cref{thm:p15-tonelli})
integrates first in \(x_k\). If \(M\) multiplies coordinate \(k\) by
\(c\neq0\), the inner integral is
\(\int g(\ldots,cx_k,\ldots)\,dx_k=|c|^{-1}\int g(\ldots,u,\ldots)\,du\) by
\cref{thm:p15-affine-lebesgue}, and \(|\det M|=|c|\). If \(M\) adds \(c\)
times coordinate \(l\neq k\) to coordinate \(k\), then with \(x_l\) held
fixed the inner integral is
\(\int g(\ldots,x_k+cx_l,\ldots)\,dx_k=\int g(\ldots,u,\ldots)\,du\) by the
same theorem with scale \(1\), and \(\det M=1\). Integrating the remaining
coordinates gives the claim in each case.

\emph{Products of good matrices are good.} If \(M_1,M_2\) are good, apply
the property of \(M_2\) to \(g\circ M_1\) and then that of \(M_1\):
\(\int g(M_1M_2x)\,dx=|\det M_2|^{-1}|\det M_1|^{-1}\int g\), and
\(\det(M_1M_2)=\det M_1\det M_2\).

\emph{Every invertible matrix is a product of elementary matrices.} Row
operations of the three kinds above are left multiplications by elementary
matrices, whose inverses are elementary of the same kind. Suppose the
invertible matrix \(N\) has its first \(j-1\) columns equal to
\(e_1,\ldots,e_{j-1}\). Some entry \(N_{ij}\) with \(i\geq j\) is nonzero;
otherwise column \(j\) would be a combination of the first \(j-1\) columns.
Exchange rows \(i\) and \(j\), scale row \(j\) by \(1/N_{ij}\), and subtract
multiples of row \(j\) from the other rows to clear column \(j\). Row \(j\)
has zeros in the first \(j-1\) columns, so those columns are unchanged.
After \(n\) steps the matrix is the identity; hence
\(E_m\cdots E_1M=I\) for elementary \(E_r\), and
\(M=E_1^{-1}\cdots E_m^{-1}\) is good.

\emph{Translations and sets.} For \(b\in\mathbb R^n\), Tonelli and
\cref{thm:p15-affine-lebesgue} applied coordinate by coordinate give
\(\int g(z+b)\,dz=\int g(z)\,dz\). Applying the good property to
\(g(\cdot+b)\) gives the displayed formula. With \(g=\mathbf 1_{MA+b}\),
which satisfies \(g(Mx+b)=\mathbf 1_A(x)\) because the map is injective and
\(MA+b\) is Borel as the inverse image of \(A\) under the continuous inverse
map, it gives \(\lambda^n(A)=|\det M|^{-1}\lambda^n(MA+b)\). The same
formula, read as \(\int g(y)\,dy=\int g(Mx+b)|\det M|\,dx\), is the
hypothesis of Definition~\ref{def:p15-w02b-p15-042-coordinate-change} with
\(j=|\det M|\).

\emph{Products.} By \cref{thm:p15-finite-product-measures},
\(\lambda^{n_1+n_2}=\lambda^{n_1}\times\lambda^{n_2}\) under the coordinate
identification. The map \(T_1\times T_2\) is a bijection of
\(U_1\times U_2\) onto \(V_1\times V_2\), Borel with Borel inverse
\(T_1^{-1}\times T_2^{-1}\). For Borel \(g\geq0\) on \(V_1\times V_2\),
Tonelli and the two hypotheses give
\[
 \iint g(u,v)\,du\,dv
  =\int_{V_2}\int_{U_1}g(T_1x,v)j_1(x)\,dx\,dv
  =\int_{U_1}j_1(x)\int_{U_2}g(T_1x,T_2y)j_2(y)\,dy\,dx ,
\]
which is the hypothesis for \(T_1\times T_2\) with factor \(j_1(x)j_2(y)\).
\end{proof}
\end{annexitem}
\begin{annexitem}{annex:p15-w02b-p15-042-diffeomorphism}{thm:p15-w02b-p15-042-diffeomorphism}{Change of variables for a \(C^1\) diffeomorphism}
\begin{proof}
Write \(S=T^{-1}:V\to U\), \(|x|_\infty=\max_i|x_i|\) for
\(x\in\mathbb R^n\), and, for a real \(n\times n\) matrix \(A\),
\(\|A\|=\max_i\sum_j|A_{ij}|\). A \emph{cube} is a set
\(C(c,r)=\{x:|x-c|_\infty\leq r\}=\prod_i[c_i-r,c_i+r]\) with \(r>0\); by
\cref{thm:p15-finite-product-measures} it has
\(\lambda^n(C(c,r))=(2r)^n\). Its boundary lies in the \(2n\) faces
\(\{x\in C(c,r):x_i=c_i\pm r\}\), each a product of intervals one of which
is a single point, so each face, and hence the boundary, is
\(\lambda^n\)-null.

\emph{(N) Norm facts.} For every \(x\),
\(|(Ax)_i|\leq\sum_j|A_{ij}||x_j|\leq\|A\|\,|x|_\infty\), so
\(|Ax|_\infty\leq\|A\|\,|x|_\infty\). Also
\(\sum_k|(AB)_{ik}|\leq\sum_j|A_{ij}|\sum_k|B_{jk}|\leq\|A\|\,\|B\|\), so
\(\|AB\|\leq\|A\|\,\|B\|\); \(\|I\|=1\); and \(\|A\|\leq n\max_{i,j}|A_{ij}|\).

\emph{Step 0: the determinant does not vanish.} Since \(S(T(x))=x\) on
\(U\) and \(T(S(y))=y\) on \(V\), the chain rule gives
\(DS(T(x))\,DT(x)=I\) and \(DT(S(y))\,DS(y)=I\). Taking determinants,
\(\det DT(x)\neq0\) and
\(|\det DT(S(y))|\,|\det DS(y)|=1\) for every \(y\in V\). The entries of
\(DT\) are continuous, hence so is \(\det DT\), and by Cramer's rule the
entries of \(DT(x)^{-1}\) are cofactors of \(DT(x)\) divided by
\(\det DT(x)\), so \(x\mapsto DT(x)^{-1}\) is continuous on \(U\).
Moreover \(T\) and \(S\) are continuous, so for Borel \(A\subseteq U\) the
set \(T(A)=S^{-1}(A)\) is Borel, and \(T\) maps open subsets of \(U\) to
open subsets of \(V\).

\emph{Step 1: a mean value bound.} Let \(F:U\to\mathbb R^n\) have
continuous first partial derivatives, let \(C=C(c,r)\subseteq U\), and
suppose \(\|DF(z)\|\leq K\) for \(z\in C\). Then
\(F(C)\subseteq\{w:|w-F(c)|_\infty\leq Kr\}\), and so
\(\lambda^n(F(C))\leq K^n\lambda^n(C)\). Indeed, for \(x\in C\) the segment
\(c+t(x-c)\), \(0\leq t\leq1\), lies in the convex set \(C\). The function
\(\phi_i(t)=F_i(c+t(x-c))\) has derivative
\(\phi_i'(t)=\sum_j\partial_jF_i(c+t(x-c))(x_j-c_j)\) by the chain rule, so
\(|\phi_i'(t)|\leq\|DF(c+t(x-c))\|\,|x-c|_\infty\leq Kr\), and the mean
value theorem gives \(|F_i(x)-F_i(c)|\leq Kr\). The set \(F(C)\) is compact,
hence Borel, and the displayed product of intervals has measure
\((2Kr)^n=K^n\lambda^n(C)\).

\emph{Step 2: closed cubes, with a continuous weight.} Let
\(h:V\to[0,\infty)\) be continuous, put
\(\Phi(x)=h(T(x))\,|\det DT(x)|\) for \(x\in U\), and let \(Q\subseteq U\)
be a cube of side \(s\). We show
\(\int_{T(Q)}h\,dy\leq\int_Q\Phi\,dx\). On the compact set \(Q\) the
continuous functions \(h\circ T\), \(|\det DT|\) and
\(x\mapsto\|DT(x)^{-1}\|\) are bounded, say by \(M_1\), \(M_2\) and \(B\),
and \(h\circ T\), \(|\det DT|\) and the entries of \(DT\) are uniformly
continuous. Fix \(0<\varepsilon\leq1\) and choose an integer \(N\geq1\) such
that for \(x,x'\in Q\) with \(|x-x'|_\infty\leq s/N\) we have
\(|h(T(x))-h(T(x'))|\leq\varepsilon\),
\(\bigl||\det DT(x)|-|\det DT(x')|\bigr|\leq\varepsilon\) and every entry of
\(DT(x)-DT(x')\) at most \(\varepsilon/(nB)\) in absolute value, so that
\(\|DT(x)-DT(x')\|\leq\varepsilon/B\) by (N).

Divide each side of \(Q\) into \(N\) closed intervals of length \(s/N\);
the products give cubes \(C_1,\ldots,C_{N^n}\) with centres \(c_k\), which
cover \(Q\), and two distinct ones meet only in boundary points. Put
\(A_k=DT(c_k)\) and \(F_k=A_k^{-1}\circ T\). Then
\(DF_k(x)=A_k^{-1}DT(x)=I+A_k^{-1}(DT(x)-A_k)\), and for \(x\in C_k\), by
(N), \(\|DF_k(x)\|\leq1+B\cdot\varepsilon/B=1+\varepsilon\). Step~1 gives
\(\lambda^n(A_k^{-1}T(C_k))\leq(1+\varepsilon)^n\lambda^n(C_k)\). The set
\(T(C_k)\) is \(A_k\) applied to the Borel set \(A_k^{-1}T(C_k)\), so
\cref{thm:p15-w02b-p15-042-linear-lebesgue} gives
\[
 \lambda^n(T(C_k))=|\det A_k|\,\lambda^n(A_k^{-1}T(C_k))
   \leq(1+\varepsilon)^n\,|\det DT(c_k)|\,\lambda^n(C_k).
\]
Let \(m_k=\max_{C_k}h\circ T\); since \(h\leq m_k\) on \(T(C_k)\),
\(\int_{T(C_k)}h\,dy\leq m_k\lambda^n(T(C_k))\). Any two points of \(C_k\)
are within \(s/N\) of each other, so for every \(x\in C_k\),
\(m_k\leq h(T(x))+\varepsilon\) and
\(|\det DT(c_k)|\leq|\det DT(x)|+\varepsilon\); multiplying and using
\(\varepsilon\leq1\),
\(m_k|\det DT(c_k)|\leq\Phi(x)+\varepsilon(M_1+M_2+1)\). Integrating over
\(C_k\) and writing \(M=M_1+M_2+1\),
\[
 \int_{T(C_k)}h\,dy\leq(1+\varepsilon)^n
   \Bigl(\int_{C_k}\Phi\,dx+\varepsilon M\lambda^n(C_k)\Bigr).
\]
Since \(T(Q)=\bigcup_kT(C_k)\), \(\int_{T(Q)}h\leq\sum_k\int_{T(C_k)}h\).
Off the null union of the boundaries,
\(\sum_k\mathbf 1_{C_k}=\mathbf 1_Q\), so \(\sum_k\int_{C_k}\Phi=\int_Q\Phi\)
and \(\sum_k\lambda^n(C_k)=\lambda^n(Q)\). Hence
\(\int_{T(Q)}h\leq(1+\varepsilon)^n(\int_Q\Phi+\varepsilon M\lambda^n(Q))\),
and letting \(\varepsilon\downarrow0\) proves the claim, the right side being
finite.

\emph{Step 3: open sets.} Let \(W\subseteq U\) be open; we show
\(\int_{T(W)}h\,dy\leq\int_W\Phi\,dx\) for \(h\) and \(\Phi\) as in Step~2.
For \(k\geq0\) let \(\mathcal D_k\) be the closed cubes
\(\prod_i[m_i2^{-k},(m_i+1)2^{-k}]\), \(m\in\mathbb Z^n\). (a) Every point
\(x\) lies in a member of \(\mathcal D_k\) (take
\(m_i=\lfloor2^kx_i\rfloor\)), and every point of that member is within
\(2^{-k}\) of \(x\) in \(|\cdot|_\infty\). (b) If \(k\leq l\),
\(Q'\in\mathcal D_k\) and \(Q\in\mathcal D_l\), then either \(Q\subseteq Q'\)
or their interiors are disjoint: in each coordinate the grid points of level
\(k\) are among those of level \(l\), so each side interval of \(Q\) lies in
one interval of level \(k\), and it lies in the corresponding side of \(Q'\)
or has interior disjoint from it; if every side lies in the corresponding
side then \(Q\subseteq Q'\), and otherwise the open interiors, which are
products of the open side intervals, are disjoint. Let \(\mathcal E_k\) be
the members of \(\mathcal D_k\) that are contained in \(W\) and in no member
of \(\mathcal E_0\cup\cdots\cup\mathcal E_{k-1}\), and
\(\mathcal E=\bigcup_k\mathcal E_k\), a countable family. Distinct members
of \(\mathcal E\) have disjoint interiors: if \(Q\in\mathcal E_l\) and
\(Q'\in\mathcal E_k\), \(k\leq l\), have meeting interiors, then
\(Q\subseteq Q'\) by (b), which for \(k<l\) contradicts the choice of
\(Q\), and for \(k=l\) forces \(Q=Q'\). The family covers \(W\): for
\(x\in W\) some sup-norm ball about \(x\) of radius \(\delta>0\) lies in
\(W\), so by (a) some member of \(\mathcal D_k\) with \(2^{-k}\leq\delta\)
contains \(x\) and lies in \(W\); for the least such \(k\), that member is in
\(\mathcal E_k\) or lies in a member of an earlier \(\mathcal E_j\), and
either way \(x\) is covered. Thus \(W=\bigcup\mathcal E\), each member is a
cube in \(U\), and off the null union of their boundaries
\(\sum_{Q\in\mathcal E}\mathbf 1_Q=\mathbf 1_W\). By Step~2,
\[
 \int_{T(W)}h\,dy\leq\sum_{Q\in\mathcal E}\int_{T(Q)}h\,dy
   \leq\sum_{Q\in\mathcal E}\int_Q\Phi\,dx=\int_W\Phi\,dx ,
\]
where the first inequality holds because \(T(W)=\bigcup_{Q\in\mathcal E}T(Q)\)
and the last equality is monotone convergence
(\cref{thm:p15-monotone-integral}) applied to the partial sums.

\emph{Step 4: equality on open sets.} Let \(W\subseteq U\) be open. Step~3
with \(h=1\) gives \(\lambda^n(T(W))\leq\int_W|\det DT|\,dx\). The inverse
\(S:V\to U\) is also a \(C^1\) diffeomorphism, \(T(W)\) is open in \(V\),
and \(|\det DT|\) is continuous and nonnegative on \(U\); Step~3 for \(S\)
with this weight gives
\[
 \int_W|\det DT(x)|\,dx=\int_{S(T(W))}|\det DT|\,dx
  \leq\int_{T(W)}|\det DT(S(y))|\,|\det DS(y)|\,dy=\lambda^n(T(W)),
\]
by Step~0. Hence \(\lambda^n(T(W))=\int_W|\det DT|\,dx\), possibly
\(\infty\).

\emph{Step 5: Borel sets.} For Borel \(A\subseteq U\) put
\(\nu(A)=\lambda^n(T(A))\) and \(\mu(A)=\int_A|\det DT|\,dx\); both are
measures on the Borel subsets of \(U\), \(\nu\) because \(T\) is a bijection
carrying Borel sets to Borel sets. Let \(d(x)\) be the sup-norm distance
from \(x\) to \(\mathbb R^n\setminus U\) (\(+\infty\) when \(U=\mathbb R^n\));
\(d\) is continuous, and \(d(x)>0\) exactly when \(x\in U\), since the
complement is closed. For \(m\geq1\) let
\(W_m=\{x:|x|_\infty<m,\ d(x)>1/m\}\). It is open, its closure lies in the
compact set \(\{x:|x|_\infty\leq m,\ d(x)\geq1/m\}\subseteq U\), on which
\(|\det DT|\) is bounded, so \(\mu(W_m)<\infty\), and
\(W_m\subseteq W_{m+1}\) with \(\bigcup_mW_m=U\). Let \(\mathcal L_m\) be the
Borel \(A\subseteq U\) with \(\nu(A\cap W_m)=\mu(A\cap W_m)\). By Step~4 it
contains every open subset of \(U\), in particular \(U\), and both measures
of \(W_m\) are finite. If \(A\subseteq B\) are in \(\mathcal L_m\), then by
item~1 of \cref{thm:p15-measure-laws}
\(\nu((B\setminus A)\cap W_m)=\nu(B\cap W_m)-\nu(A\cap W_m)\), and the same
for \(\mu\); increasing unions are handled by item~3 of that theorem. So
\(\mathcal L_m\) is a lambda-system on \(U\). The open subsets of \(U\) form
a pi-system, and they generate the Borel subsets of \(U\) because \(U\) is
open. By \cref{thm:p15-measure-uniqueness}, \(\mathcal L_m\) contains every
Borel \(A\subseteq U\). Letting \(m\to\infty\) and using item~3 again,
\(\nu(A)=\mu(A)\), that is, \(\lambda^n(T(A))=\int_A|\det DT|\,dx\).

\emph{Step 6: functions.} For Borel \(E\subseteq V\), take \(A=S(E)\),
Borel by Step~0 applied to \(S\); then \(T(A)=E\) and
\(\mathbf 1_A=\mathbf 1_E\circ T\), so Step~5 reads
\(\int_V\mathbf 1_E\,dy=\int_U\mathbf 1_E(T(x))|\det DT(x)|\,dx\).
Both sides are linear in \(g\), so the identity holds for nonnegative simple
\(g\). A Borel \(g:V\to[0,\infty]\) is the increasing limit of nonnegative
simple Borel \(g_k\) (\cref{thm:p15-measurable-limits}), and
\(g_k\circ T\) increases to \(g\circ T\); monotone convergence
(\cref{thm:p15-monotone-integral}) on both sides gives the identity for
\(g\). Finally \(T\) and \(S\) are continuous, hence Borel, and
\(|\det DT|\) is a continuous, hence Borel, function into \((0,\infty)\), so
\((T,|\det DT|)\) satisfies Definition~\ref{def:p15-w02b-p15-042-coordinate-change}.
\end{proof}
\end{annexitem}

\begin{annexitem}{annex:p15-w02b-p15-042-coordinate-change-thm}{thm:p15-w02b-p15-042-coordinate-change-thm}{Item 8: entropy in new coordinates}
\begin{proof}
The function \(p_y\) is Borel, since \(p\), \(j\) and \(T^{-1}\) are, and
\(p_y(T(x))=p(x)/j(x)\) on \(U\). For Borel \(B\subseteq\mathbb R^n\), apply
the hypothesis to \(g=\mathbf 1_{B\cap V}\,p_y\):
\[
 \int_Bp_y\,dt=\int_U\mathbf 1_B(T(x))\frac{p(x)}{j(x)}j(x)\,dx
              =P(x\in U,\ T(x)\in B)=P(y\in B).
\]
So \(p_y\) is a density of \(y\). Apply the hypothesis to
\(g=\eta(p_y)^\pm\) on \(V\); since \(j>0\),
\[
 \int\eta(p_y)^\pm\,dt=\int_U\bigl(j\,\eta(p/j)\bigr)^\pm\,dx ,
\]
and \(\eta(p_y)=0\) off \(V\). Where \(p>0\),
\(j\eta(p/j)=\eta(p)+p\log j\); where \(p=0\) both sides vanish. Because
\(p=0\) a.e.\ off \(U\), the integrals over \(U\) are integrals over
\(\mathbb R^n\). Hence \(H(y)\) is defined exactly when
\(\eta(p)+p\log j\) is quasi-integrable, and then equals its integral. Under
either hypothesis (E2) gives quasi-integrability and
\(H(y)=H(x)+\int p\log j\,dx\). Since \(J(x/y)=1/j(x)\) at \(y=T(x)\),
\(\log j=-\log J(x/y)\), which is the second form.
\end{proof}
\end{annexitem}

\begin{annexitem}{annex:p15-w02b-p15-042-linear-change}{thm:p15-w02b-p15-042-linear-change}{Item 9: linear coordinates}
\begin{proof}
In matrix form \(y=Mx\) with \(M_{ji}=a_{ij}\), so \(M\) is the transpose of
\((a_{ij})\) and \(|\det M|=|\det(a_{ij})|\). By
\cref{thm:p15-w02b-p15-042-linear-lebesgue} this is a coordinate change on
\(\mathbb R^n\) with the constant factor \(j=|\det(a_{ij})|\), and
\(\int p|\log j|\,dx=|\log j|<\infty\). \Cref{thm:p15-w02b-p15-042-coordinate-change-thm}
gives \(H(y)=H(x)+\log|\det(a_{ij})|\). A rotation has
\(|\det M|=1\); a translation has factor \(1\) by the same lemma; and for any
coordinate change with \(j\equiv1\) the added term is zero.
\end{proof}
\end{annexitem}

\begin{annexitem}{annex:p15-w02b-p15-042-difference-invariance}{thm:p15-w02b-p15-042-difference-invariance}{Differences that do not depend on coordinates}
\begin{proof}
By \cref{thm:p15-w02b-p15-042-coordinate-change-thm},
\(H(x')=H(x)+\int p(x)\log j_1\,dx\) and
\(H(y')=H(y)+\int p(y)\log j_2\,dy\), both finite. By
\cref{thm:p15-w02b-p15-042-linear-lebesgue}, \(T_1\times T_2\) is a coordinate
change with factor \(j_1(x)j_2(y)\), and \(P((x,y)\in U_1\times U_2)=1\). By
Tonelli,
\(\iint p(x,y)|\log j_1(x)+\log j_2(y)|\leq\int p(x)|\log j_1|+\int p(y)|\log j_2|<\infty\),
and by Fubini
\(\iint p(x,y)\log(j_1(x)j_2(y))=\int p(x)\log j_1+\int p(y)\log j_2\).
Hence \(H(x',y')=H(x,y)+\int p(x)\log j_1+\int p(y)\log j_2\), finite. The
added terms cancel in \(H(x')+H(y')-H(x',y')\). Finally
\(H(x)+H(y)-H(x,y)=H(x)-H_y(x)\) by \cref{thm:p15-w02b-p15-042-conditional}.
\end{proof}
\end{annexitem}

\begin{annexitem}{annex:p15-w02b-p15-042-boundaries}{thm:p15-w02b-p15-042-boundaries}{Coordinate-relative values, rankings and divergence}
\begin{proof}
The map \(T(x)=x^2\) is an increasing \(C^1\) bijection of \((0,1)\) onto
itself with \(T'(x)=2x>0\), so by \cref{thm:p15-criteria-real-integration}
it is a coordinate change with \(j(x)=2x\). Two integrals are needed. First,
\(\int_{(0,1)}\log x\,dx=-1\): on \([\varepsilon,1]\) the integral of the
continuous function \(-\log x\geq0\) is
\(1+\varepsilon\log\varepsilon-\varepsilon\) by the fundamental theorem of
calculus, and monotone convergence lets \(\varepsilon\downarrow0\), since
\(\varepsilon\log\varepsilon\to0\). Second,
\(\int_{(0,1)}x\log x\,dx=-\tfrac14\): the integrand is continuous and
bounded on \((0,1]\), with antiderivative \(\tfrac12x^2\log x-\tfrac14x^2\)
tending to \(0\) at \(0\), and dominated convergence applies on
\([\varepsilon,1]\).

\emph{Item 1.} \(H(x)=0\) by the equality case of
\cref{thm:p15-w02b-p15-042-volume} with \(v=1\). Since
\(\int|\log(2x)|\,dx<\infty\),
\cref{thm:p15-w02b-p15-042-coordinate-change-thm} gives
\(H(y)=0+\log2+\int\log x\,dx=\log2-1\).

\emph{Item 2.} \(P(y\leq c)=P(x\leq\sqrt c)=\int_0^{\sqrt c}2x\,dx=c\) for
\(0<c<1\), so \(y\) is uniform on \((0,1)\) (the intervals \((0,c]\) form a
pi-system generating the Borel sets of \((0,1)\);
\cref{thm:p15-measure-uniqueness}), and \(H(y)=0\). Directly,
\(H(x)=-\int2x(\log2+\log x)\,dx=-\log2+\tfrac12\). The theorem confirms
this: \(\int2x|\log(2x)|\,dx<\infty\) and
\(H(y)=\tfrac12-\log2+\log2+\int2x\log x\,dx=0\).

\emph{Item 3.} The map \(T(x)=e^{-1/x}\) is \(C^1\) on \((0,1)\) with
\(T'(x)=e^{-1/x}/x^2>0\), \(T(x)\to0\) as \(x\downarrow0\) and
\(T(x)\to e^{-1}\) as \(x\uparrow1\); so it is an increasing bijection onto
\((0,e^{-1})\) and a coordinate change with \(j=T'\). Here
\(\log j(x)=-1/x-2\log x\). Its maximum on \((0,1)\) occurs where
\(1/x^2-2/x=0\), at \(x=\tfrac12\), with value \(-2+2\log2<0\); so
\(\log j<0\), and
\(\int\log j\,dx=-\int_{(0,1)}x^{-1}\,dx-2\int\log x\,dx=-\infty\), the first
integral being infinite by the substitution \(u=1/x\) and the computation of
\(\int_{(1,\infty)}u^{-1}\,du\) in the proof of
\cref{thm:p15-w02b-p15-042-defined-cases}. Since \(H(x)=0\) is finite and
\(\int p\log j\) exists with value \(-\infty\), the second hypothesis of
\cref{thm:p15-w02b-p15-042-coordinate-change-thm} gives \(H(y)=-\infty\).

\emph{Item 4.} By \cref{thm:p15-normal-sums}, \(x_1=\sigma_1z_1\) has law
\(\operatorname{normal}(0,\sigma_1^2)\), and \(x_2\), the sum of the
independent variables \(\sigma_2\rho z_1\) and \(\sigma_2\sqrt{1-\rho^2}z_2\),
has law \(\operatorname{normal}(0,\sigma_2^2\rho^2+\sigma_2^2(1-\rho^2))\).
So \(H(x_i)=\log\sqrt{2\pi e}\,\sigma_i\) by
\cref{thm:p15-w02b-p15-042-gaussian}, item~5. The pair \(z=(z_1,z_2)\) has
density \(g\) of the Gaussian proof with \(n=2\), and \(H(z)=\log2\pi e\) by
the same item with \(a\) the identity. Further \(x=Mz\) with
\(M=\begin{pmatrix}\sigma_1&0\\ \sigma_2\rho&\sigma_2\sqrt{1-\rho^2}\end{pmatrix}\),
\(\det M=\sigma_1\sigma_2\sqrt{1-\rho^2}>0\). By
\cref{thm:p15-w02b-p15-042-linear-change},
\(H(x_1,x_2)=\log2\pi e+\log(\sigma_1\sigma_2\sqrt{1-\rho^2})\), and
\[
 H(x_1)+H(x_2)-H(x_1,x_2)
  =\log(2\pi e\,\sigma_1\sigma_2)-\log(2\pi e\,\sigma_1\sigma_2\sqrt{1-\rho^2})
  =-\tfrac12\log(1-\rho^2).
\]
\end{proof}
\end{annexitem}

\begin{annexitem}{annex:p15-w02b-p15-042-power-coordinates}{thm:p15-w02b-p15-042-power-coordinates}{Uniform laws in power coordinates}
\begin{proof}
Put \(t=\ell^a\). The map \(s\mapsto s^a\) is strictly increasing on
\((0,1)\), so \(\{\ell\leq c\}=\{t\leq c^a\}\) and \(P(\ell\leq c)=c^a\)
for \(0<c<1\). In particular different exponents \(a\) give different
laws of \(\ell\), since \(P(\ell\leq\tfrac12)=2^{-a}\).

Since \(\ell^b=(\ell^a)^{b/a}\), we have \(\ell^b=T(t)\) with
\(T(t)=t^r\). The map \(T\) is \(C^1\) on \((0,1)\) with
\(T'(t)=rt^{r-1}>0\), \(T(t)\to0\) as \(t\downarrow0\) and \(T(t)\to1\) as
\(t\uparrow1\); so it is an increasing bijection of \((0,1)\) onto itself,
and by item~2 of
Definition~\ref{def:p15-w02b-p15-042-coordinate-change} it is a coordinate
change with factor \(j=T'\).

The variable \(t\) has Lebesgue density \(1\) on \((0,1)\), so \(H(t)=0\) by
the equality case of \cref{thm:p15-w02b-p15-042-volume} with
\(\mu=\lambda^1\), \(A=(0,1)\) and \(v=1\). Next,
\(\log j(t)=\log r+(r-1)\log t\), and
\[
 \int_{(0,1)}|\log j(t)|\,dt\leq|\log r|+|r-1|\int_{(0,1)}|\log t|\,dt
   =|\log r|+|r-1|<\infty ,
\]
using \(\int_{(0,1)}\log t\,dt=-1\), proved in the annex item for
\cref{thm:p15-w02b-p15-042-boundaries}. The first set of hypotheses of
\cref{thm:p15-w02b-p15-042-coordinate-change-thm} therefore holds, and
\[
 H(\ell^b)=H(t)+\int_{(0,1)}\bigl(\log r+(r-1)\log t\bigr)\,dt
          =\log r-(r-1)=\log r+1-r .
\]

For the inequality let \(f(r)=\log r+1-r\) on \((0,\infty)\). It is
differentiable with \(f'(r)=1/r-1\), which is positive on \((0,1)\) and
negative on \((1,\infty)\). By the mean value theorem \(f\) is strictly
increasing on \((0,1]\) and strictly decreasing on \([1,\infty)\), so
\(f(r)<f(1)=0\) for every \(r\neq1\). As \(r=b/a\), equality holds if and only
if \(a=b\).

For the final sentence, fix \(b\). The law for which \(\ell^b\) is uniform
has \(H(\ell^b)=f(1)=0\). Any law for which \(\ell^a\) is uniform with
\(a\neq b\) is a different law, by the first paragraph, and has
\(H(\ell^b)=f(b/a)<0\).
\end{proof}
\end{annexitem}
